\documentclass[11pt]{article}
\usepackage[margin=1in]{geometry}
\usepackage{amsmath,amssymb,amsthm}
\newtheorem{theorem}{Theorem}[section]
\newtheorem{proposition}[theorem]{Proposition}
\title{Crossing-Point Analysis for Weighted Rolling Validation: Theory}
\author{Aileen Tao}
\date{}
\begin{document}
\maketitle

\section{Question and approach}

Rolling validation (RV) selects between models by adding their validation losses over successive training sizes. The target, however, is the model with smaller population prediction risk at the \emph{current} training size. If that preferred model changes as the training sample grows, the two comparisons need not change sign together: RV still includes losses from earlier, smaller training samples. The question of this report is whether weighting later validation observations can bring the expected RV crossing close to the population-risk crossing.

The two tasks are distinct. First, expected RV should change preference at approximately the right training size (\emph{alignment}). Second, the realized, random RV comparison should be sufficiently concentrated around its mean (\emph{concentration}). Weighting can improve the first task while making the second harder.

\subsection*{Why Analysis 08 is not enough}

Analysis 08 considered a sample mean against an oracle that was always better in population risk. It studied the mean and variance of the RV loss difference and whether their ratio gave a selection-consistency bound. There was no preference reversal, hence no crossing to align. The present model deliberately has a reversal. Its nonzero mean parameter also creates a new linear random term in RV, so its variance cannot be imported from 08 without calculation.

Analysis 09 therefore introduces a crossing model in which the population-preferred
model changes with training size. The report first studies the population crossing
and the later ordinary-RV crossing. It then asks whether weighting can align the
expected RV crossing while retaining enough concentration for reliable selection.

\section{The two comparisons studied in 09}

For models 0 and 1 trained on $n$ observations, write
\[
\Delta_R(n)=R_0(n)-R_1(n).
\]
The population-risk difference $\Delta_R(n)$ compares the two models only at the
\emph{current} training size $n$. It is not accumulated over previous training
sizes. By contrast, ordinary RV accumulates validation-loss differences from
earlier training sizes. Its expected difference $\Delta_{\rm RV}(n)$ therefore
contains historical contributions even after current population risk has
changed its preferred model.

A negative value favors Model 0, a positive value favors Model 1, and zero means equal population risks. If $RV_{k,w}(n)$ denotes Model $k$'s weighted RV loss through the validation time with $n$ training observations, write
\[
\Delta_w(n)=\mathbb E[RV_{0,w}(n)-RV_{1,w}(n)].
\]
Its sign describes the preference of \emph{expected} weighted RV, not necessarily a particular realized selection. Thus $\Delta_R(n)$ tells us what the population comparison prefers, while $\Delta_w(n)$ tells us what expected weighted RV prefers. We compare their first strict negative indices:
\[
n_{\rm pop}=\min\{n\ge1:\Delta_R(n)<0\},\qquad
n_{{\rm RV},w}=\min\{n\ge1:\Delta_w(n)<0\}.
\]
The sets are shown to be nonempty below for the weights used here. The goal is to make these indices close, principally in the ratio sense as the crossing moves outward.

\section{A model with a genuine crossing}

Let
\[
Y_i=\theta+\xi_i,\qquad \xi_i\overset{\rm iid}{\sim}N(0,1),
\qquad 0<|\theta|<1.
\]
Model 0 predicts a new response using the training sample mean $\bar Y_n$; Model 1 always predicts zero. Both use squared loss. Model 0 is unbiased but has estimation variance. Model 1 has no estimation variance but retains squared bias $\theta^2$. Hence the zero predictor can be better early, and the sample mean better later. A fixed shrinkage predictor would give a similar bias--variance crossing, but the zero predictor keeps the mechanism and formulas transparent.

\subsection{Population prediction risk}

Here $R_k(n)$ averages squared prediction loss over the training data and an
independent new response. The exact crossing is as follows.

\begin{theorem}[Population crossing]\label{thm:population}
\[
R_0(n)=1+\frac1n,\qquad R_1(n)=1+\theta^2,\qquad
\Delta_R(n)=\frac1n-\theta^2,
\qquad n_{\rm pop}=\lfloor\theta^{-2}\rfloor+1.
\]
\end{theorem}
\noindent\textbf{Interpretation.} The common $1$ is new-response noise.
The comparison is between the sample mean's estimation variance and the
zero predictor's squared bias. The displayed index is the first \emph{strict}
Model-0 preference; an equality at $n=\theta^{-2}$ is a tie.

\noindent\textbf{Proof.} See Appendix~\ref{app:pop}.

\subsection{Why ordinary RV crosses late}

At validation time $m+1$, Model 0 has been trained on $m$ observations. Its
expected loss difference relative to Model 1 is $1/m-\theta^2$. Ordinary RV
accumulates these one-time differences:
\begin{equation}\label{eq:ordinary}
\Delta_{\rm RV}(n)=\sum_{m=1}^n\left(\frac1m-\theta^2\right).
\end{equation}
Define
\[
H_n=\sum_{m=1}^n\frac1m=1+\frac12+\frac13+\cdots+\frac1n.
\]
$H_n$ is the $n$-th harmonic number, and
$\Delta_{\rm RV}(n)=H_n-n\theta^2$. Since $H_n\sim\log n$, positive
contributions from early validation times accumulate approximately
logarithmically. The one-time expectation and sum are derived in
Appendix~\ref{app:ordinary-mean}.

The distinction is temporal: $\Delta_R(n)$ compares population risks only
at the current training size $n$, whereas $\Delta_{\rm RV}(n)$ is an
accumulated expected comparison. After the current population risk begins
to favor Model 0, ordinary RV can still favor Model 1 because its sum
retains positive contributions from earlier, smaller training sizes.

\begin{proposition}[Ordinary-RV lag]\label{prop:ordinary}
As $\theta\to0$,
\[
n_{\rm RV}\sim\theta^{-2}\log(\theta^{-2}),\qquad
\frac{n_{\rm RV}}{n_{\rm pop}}\sim\log(\theta^{-2})\to\infty.
\]
\end{proposition}
\noindent\textbf{Interpretation.} The accumulated history delays the
expected ordinary-RV crossing by a growing factor relative to the
population crossing.

\noindent\textbf{Proof.} See Appendix~\ref{app:ordinary-cross}.

\section{Weighted RV and crossing alignment}

The lag originates in accumulated early contributions. Increasing weights
give more influence to later validations, which use larger training samples.
For deterministic $w_i>0$, independent of the data, define the realized
weighted-RV loss difference and its expectation by
\[
\Xi_{n,w}=\sum_{i=2}^{n+1}w_i\{(Y_i-\bar Y_{i-1})^2-Y_i^2\},
\qquad \Delta_w(n)=\mathbb E[\Xi_{n,w}].
\]
A negative value selects Model 0. The final validation prediction uses
$n$ training observations. Multiplying all weights by one positive
constant changes neither sign nor crossing.

At validation time $i$, the expected loss difference is
\[
d_i=\frac{1}{i-1}-\theta^2.
\]
Define the weighted historical reciprocal sum and the total weight:
\[
A_n=\sum_{i=2}^{n+1}\frac{w_i}{i-1},\qquad
W_n=\sum_{i=2}^{n+1}w_i.
\]
The expected weighted-RV difference is
\begin{equation}\label{eq:weighted-mean}
\Delta_w(n)=A_n-\theta^2W_n.
\end{equation}
Here $A_n$ collects the estimation-variance terms from all validation
times, while $\theta^2W_n$ collects the bias terms. The derivation of
$d_i$ and \eqref{eq:weighted-mean} is in
Appendix~\ref{app:weighted-mean}.

Normalizing by $W_n$ makes $A_n/W_n$ the weighted average of the
historical reciprocal training sizes. Population risk uses only the
current value $1/n$. Their exact relationship is
\begin{equation}\label{eq:remainder}
\frac{\Delta_w(n)}{W_n}
=\Delta_R(n)+\left(\frac{A_n}{W_n}-\frac1n\right).
\end{equation}
The correction in parentheses is the effect of earlier training sizes.
It is nonnegative, and strictly positive for $n>1$ with positive
weights, so the expected weighted-RV crossing cannot precede the
population crossing. The identity and this crossing order are proved
in Appendix~\ref{app:remainder}; crossing monotonicity and existence
are established in Appendix~\ref{app:monotonicity}.

The alignment question is visible by comparing
\[
\frac{\Delta_w(n)}{W_n}=\frac{A_n}{W_n}-\theta^2,
\qquad
\Delta_R(n)=\frac1n-\theta^2.
\]
Their only difference is $A_n/W_n$ versus $1/n$. To make the expected
weighted comparison resemble the current population comparison near
their moving crossing, the target is $A_n/W_n\sim1/n$, equivalently
$nA_n/W_n\to1$. The next theorem gives this condition a precise
crossing interpretation.

\begin{theorem}[Sufficient crossing alignment]\label{thm:alignment}
If the positive deterministic weights do not depend on $\theta$ and
\begin{equation}\label{eq:alignment}
\frac{nA_n}{W_n}\longrightarrow1,
\end{equation}
then $n_{{\rm RV},w}$ exists for all sufficiently small nonzero $\theta$ and
\[
\frac{n_{{\rm RV},w}}{n_{\rm pop}}\longrightarrow1
\qquad(\theta\to0).
\]
\end{theorem}
\noindent\textbf{Interpretation.} If the weighted average of reciprocal
training sizes approaches the latest reciprocal training size, the
\emph{expected} weighted-RV crossing aligns with the population
crossing. The theorem does not control random fluctuations.

\noindent\textbf{Proof.} See Appendix~\ref{app:alignment}.

\subsection{Representative weights}

Crossing alignment requires $A_n/W_n\sim1/n$, equivalently
$nA_n/W_n\to1$. The following families show how different rates of
weight growth affect this condition.

\paragraph{Log-power weights.}
Consider $w_i=[\log(i+1)]^p$ for $p>0$. These weights test whether a
gradual emphasis on later observations can remove the lag. Appendix
\ref{app:logpower} establishes
\[
\frac{A_n}{W_n}\sim\frac{\log n}{(p+1)n},
\qquad
\frac{n_{{\rm RV},w}}{n_{\rm pop}}
\sim\frac{\log(\theta^{-2})}{p+1}.
\]
The target is $1/n$, but the additional factor $\log n/(p+1)$
diverges. These weights therefore fail alignment, and their crossing
ratio diverges rather than tending to $1$. The case $p=0$ is ordinary
unweighted RV, $p=1$ is the log weight, and $p=2$ is the squared-log
weight. Increasing $p$ reduces the finite crossing lag without
removing its asymptotic divergence. The full asymptotic derivation is
in Appendix~\ref{app:logpower}.

\paragraph{Log-log weight.}
The slower-growing choice $w_i=\log\log(i+e)$ checks whether a gentler
increase behaves differently. It satisfies
\[
\frac{A_n}{W_n}\sim\frac{\log n}{n}.
\]
This still exceeds the $1/n$ target by a factor of $\log n$.
Although it grows much more slowly than log-power weights, it also
fails asymptotic alignment: its crossing has the same leading lag
ratio as ordinary RV. The derivation is in Appendix~\ref{app:loglog}.

\paragraph{Geometric weights.}
Consider $w_i=q^i$ for $q>1$. These weights grow much faster, so most
of the total weight falls on the latest validation observations. Their
training sizes are close to $n$, making their reciprocal training
sizes close to $1/n$. Appendix~\ref{app:geometric} proves
\[
\frac{nA_n}{W_n}\to1.
\]
Thus geometric weights satisfy Theorem~\ref{thm:alignment}, and
$n_{{\rm RV},w}/n_{\rm pop}\to1$: the \emph{expected} weighted-RV
crossing aligns asymptotically with the population crossing. The
same appendix calculates
\[
\frac{\sum_{i=2}^{n+1}w_i^2}{W_n^2}
\to\frac{q-1}{q+1}>0.
\]
The geometric weights are highly concentrated on a small number of
late validation times, so they achieve crossing alignment but retain
relatively little averaging across validation observations.

The representative weights reveal a contrast. Slowly increasing
log-type weights spread weight across many validation times and
reduce finite lags, but their historical average $A_n/W_n$ remains
too large relative to $1/n$ for asymptotic alignment. Geometric
weights align the crossing by concentrating on the latest times,
which reduces averaging. This motivates the next section.

\section{Randomness near the moving crossing}

The population crossing satisfies $n_{\rm pop}\sim\theta^{-2}$ as
$\theta\to0$, so $n/n_{\rm pop}$ is approximately $n\theta^2$.
The quantity $n\theta^2$ tells us where the current training size
lies relative to the moving population crossing. Letting
$\theta\to0$ sends that crossing to larger sample sizes, while
requiring
\[
n\theta^2\to c\in(0,\infty)
\]
keeps $n$ at a fixed relative position. Here $c<1$ is before the
population crossing, $c=1$ is at its leading-order location, and
$c>1$ is after it. This is the relevant regime because the research
question asks whether weighted RV behaves correctly near the point
where the population-preferred model changes. The precise sign
calculation is in Appendix~\ref{app:scaling}.

Crossing alignment concerns the expectation of weighted RV. The
realized statistic $\Xi_{n,w}$ is random, so its variance also matters
near the crossing. A sufficient Chebyshev concentration condition,
when its mean is nonzero, is
\[
\frac{\operatorname{Var}(\Xi_{n,w})}
{\{\mathbb E\Xi_{n,w}\}^2}\to0.
\]
This would control the probability of a sign error relative to the
expected statistic.

The 08 Gaussian toy model used $\theta=0$. Here
$Y_i=\theta+\xi_i$ with $\theta\ne0$. Expanding a validation loss
difference produces an additional linear Gaussian term
$-2\theta\xi_i$. It vanishes when $\theta=0$, which is why it was
absent in 08. After weighting and summing, its random contribution
is $-2\theta\sum_{i=2}^{n+1}w_i\xi_i$, with variance
$4\theta^2\sum_{i=2}^{n+1}w_i^2$. Concentrating weights on the latest
times can improve alignment, but also makes this noise large relative
to the squared total weight.

The exact variance identity is
\begin{equation}\label{eq:variance}
\operatorname{Var}(\Xi_{n,w})
=2\sum_{t=1}^{n}a_t^2+
4\sum_{l=2}^{n+1}(l-1)b_l^2+
4\theta^2\sum_{i=2}^{n+1}w_i^2,
\end{equation}
where $a_t=\sum_{m=t}^nw_{m+1}/m^2$ for $1\le t\le n$ and
$b_l=\sum_{m=l}^nw_{m+1}/m^2-w_l/(l-1)$ for $2\le l\le n+1$,
with empty sums set to zero. The first two terms describe quadratic
Gaussian noise; the last is the new linear noise. The loss expansion
and variance proof are entirely in Appendices~\ref{app:loss-expansion}
and~\ref{app:full-var}.

Crossing alignment requires increasing emphasis on the latest
validation times, since only their reciprocal training sizes are
close to the current population term $1/n$. Concentrating weight
near the end reduces averaging across validation observations. The
expected crossing may therefore align while the realized statistic
remains too noisy near it. To describe the amount of averaging,
consider
\[
\frac{\sum_{i=2}^{n+1}w_i^2}{W_n^2}.
\]
This measures concentration of the normalized weights: it is small
when weight is spread across many observations and larger when only
a few observations dominate. Its reciprocal is an effective number
of observations. Thus divergence of
$n\sum_iw_i^2/W_n^2$ means the effective amount of averaging is
much smaller than order $n$.

\begin{theorem}[Near-crossing obstruction]\label{thm:tradeoff}
Under \eqref{eq:alignment},
\[
n\frac{\sum_{i=2}^{n+1}w_i^2}{W_n^2}\longrightarrow\infty.
\]
For every sequence $\theta\to0$ and $n\theta^2\to c\in(0,\infty)\setminus\{1\}$,
\begin{equation}\label{eq:obstruction}
\frac{\operatorname{Var}(\Xi_{n,w})}
{\{\mathbb E\Xi_{n,w}\}^2}\longrightarrow\infty.
\end{equation}
\end{theorem}
\noindent\textbf{Interpretation.} Near the moving crossing,
alignment and Chebyshev concentration cannot hold simultaneously in
this 09 model. Strong late weighting reduces averaging and makes the
new linear Gaussian noise too large relative to the squared mean.
The Chebyshev ratio therefore diverges rather than vanishes. This
criterion does not determine the limiting wrong-selection probability.
The complete proof and variance calculation are in
Appendices~\ref{app:effective}--\ref{app:obstruction}
and~\ref{app:loss-expansion}--\ref{app:full-var}.

\paragraph{Logarithmic and log-log weights.}
These spread mass across more validation times and retain better
averaging, but their expected crossings remain asymptotically late.
For $c>1$, they can concentrate around an expected sign that
disagrees with current population preference.

\paragraph{Geometric weights.}
These align the expected crossing by concentrating strongly on the
latest validation times, but the same concentration causes the
near-crossing Chebyshev ratio to diverge.

For fixed $\theta\ne0$ and $n\to\infty$, some intermediate-growth
weights can achieve both alignment and concentration, but this is a
different and weaker regime than the moving-crossing problem.

\section{What 09 establishes}

Analysis 09 shows that ordinary rolling validation changes model
preference later than population risk because it accumulates
information from earlier training sizes. Weighting can reduce this
lag, and sufficiently concentrated late-stage weights can align the
expected crossing with the population crossing. Near the moving
crossing, however, such concentration reduces averaging and
prevents the Chebyshev concentration condition from holding.
In this model, the same positive deterministic weighting behavior
cannot achieve both properties.

\section{What remains open}

The result is specific to the 09 model, positive deterministic
weights, and the near-crossing Chebyshev argument. It does not
determine the exact limiting wrong-selection probability. Future
work could study that probability directly or examine other
weighting schemes, validation designs, or crossing models.

\appendix
\section{Population risk and ordinary RV}

\subsection{Population-risk formulas and strict crossing}\label{app:pop}
Let $Y_{\rm new}=\theta+\xi_{\rm new}$ be independent of the training sample. Since $\bar Y_n=\theta+S_n/n$ with $S_n=\sum_{t=1}^n\xi_t$, the Model-0 error is $\xi_{\rm new}-S_n/n$. Independence, zero means, and $\operatorname{Var}(S_n)=n$ give $R_0(n)=1+1/n$. Model 1 has error $\theta+\xi_{\rm new}$, hence $R_1(n)=1+\theta^2$. Their difference is $1/n-\theta^2$, negative exactly when $n>\theta^{-2}$. The least integer strictly above $\theta^{-2}$ is $\lfloor\theta^{-2}\rfloor+1$, including when $\theta^{-2}$ itself is an integer (the equality index is a tie). This proves Theorem~\ref{thm:population}.

\subsection{One-time and ordinary-RV expectation}\label{app:ordinary-mean}
At time $m+1$, $\bar Y_m=\theta+S_m/m$ and the fresh error of Model 0 is $\xi_{m+1}-S_m/m$; that of Model 1 is $\theta+\xi_{m+1}$. Expanding and using independence of $\xi_{m+1}$ and $S_m$ gives
\[
\mathbb E[(Y_{m+1}-\bar Y_m)^2-Y_{m+1}^2]
=(1+1/m)-(1+\theta^2)=1/m-\theta^2.
\]
Summing over $m=1,\ldots,n$ proves \eqref{eq:ordinary}.

\subsection{Ordinary-RV crossing asymptotics}\label{app:ordinary-cross}
Put $h=\theta^2$ and $N=h^{-1}$. The mean $H_n-nh$ has increments $1/(n+1)-h$, so it is initially increasing and eventually decreasing to $-\infty$; its first strict negative index exists and is unique. Integral bounds give $\log(n+1)\le H_n\le1+\log n$. Fix $\varepsilon\in(0,1)$ and set
\[
n_-=\lfloor(1-\varepsilon)N\log N\rfloor,\qquad
n_+=\lceil(1+\varepsilon)N\log N\rceil.
\]
Both tend to infinity; $\log n_\pm=\log N+\log\log N+O(1)$, while $n_\pm h=(1\pm\varepsilon)\log N+o(1)$. Thus $H_{n_-}-n_-h=\varepsilon\log N+o(\log N)>0$ and $H_{n_+}-n_+h=-\varepsilon\log N+o(\log N)<0$. The crossing lies between $n_-$ and $n_+$. Divide by $N\log N$ and let $\varepsilon\downarrow0$ to obtain $n_{\rm RV}\sim N\log N$. Since $n_{\rm pop}=N+O(1)$, Proposition~\ref{prop:ordinary} follows.

\section{Weighted RV and crossing alignment}

\subsection{Weighted expectation}\label{app:weighted-mean}
Apply the one-time expectation in Appendix~\ref{app:ordinary-mean} with $m=i-1$. Deterministic weights and linearity of expectation give
\[
\Delta_w(n)=\sum_{i=2}^{n+1}w_i\bigl((i-1)^{-1}-\theta^2\bigr)
=\sum_{i=2}^{n+1}\frac{w_i}{i-1}
-\theta^2\sum_{i=2}^{n+1}w_i=A_n-\theta^2W_n.
\]
No independence between validation losses is asserted or needed.

\subsection{Exact remainder and crossing order}\label{app:remainder}
Divide the preceding identity by $W_n>0$ and add and subtract $1/n$ to obtain \eqref{eq:remainder}. Each $1/(i-1)\ge1/n$. Their positive-weighted average therefore satisfies $A_n/W_n\ge1/n$; for $n>1$, at least one earlier term is strictly larger and has positive weight, so the inequality is strict. If $\Delta_R(n)\ge0$, then $\Delta_w(n)/W_n\ge0$. Consequently a strict weighted-RV crossing cannot precede $n_{\rm pop}$. Equality is possible at $n=1$ or if nonnegative (rather than positive) weights are supported only on the final term.

\subsection{Monotonicity and crossing existence}\label{app:monotonicity}
Write $r_n=A_n/W_n$. Adding the next validation time gives
\[
r_{n+1}=\frac{W_nr_n+w_{n+2}/(n+1)}{W_n+w_{n+2}}.
\]
Every term in $r_n$ is at least $1/n>1/(n+1)$, so $r_{n+1}<r_n$ for positive weights. Thus the equation $r_n<h$ has at most one first strict crossing. It exists exactly when $\inf_n r_n<h$; since $r_n$ decreases, this is equivalent to $\lim_n r_n<h$. If the limit equals $h$, strict crossing need not occur. Under \eqref{eq:alignment}, $r_n\sim1/n\to0$, so crossing exists for every fixed $h>0$; $r_1=1>h$ because $0<h<1$.

\subsection{Sufficient-alignment theorem}\label{app:alignment}
Let $h=\theta^2\downarrow0$ and fix $\varepsilon\in(0,1)$. Set $n_-=\lfloor(1-\varepsilon)/h\rfloor$ and $n_+=\lceil(1+\varepsilon)/h\rceil$. Both grow, and \eqref{eq:alignment} gives $r_{n_\pm}=(1+o(1))/n_\pm$. Therefore
\[
\frac{r_{n_-}}h\to\frac1{1-\varepsilon}>1,\qquad
\frac{r_{n_+}}h\to\frac1{1+\varepsilon}<1.
\]
By Appendix~\ref{app:monotonicity}, the first strict crossing lies in $(n_-,n_+]$ for small $h$. Rounding changes either endpoint by at most one. Divide by $h^{-1}$ and let $\varepsilon\downarrow0$; since $n_{\rm pop}=h^{-1}+O(1)$, the crossing ratio tends to one. This proves Theorem~\ref{thm:alignment}.

\subsection{Log-power weights and their crossings}\label{app:logpower}
Let $w_i=[\log(i+1)]^p$ with fixed $p\ge0$. A shift by one changes $\log(i+1)$ by $O(1/i)$, so its $p$th power has the same leading order as $(\log i)^p$ (with $p=0$ interpreted as $1$). Integral comparison, or summation by parts, yields
\[
A_n\sim\int_2^n\frac{(\log x)^p}{x}\,dx
=\frac{(\log n)^{p+1}}{p+1}.
\]
For completeness, the summand is eventually decreasing: its logarithmic derivative is $(p/\log x-1)/x<0$ for $x>e^p$. Thus its sum differs from the integral by a bounded initial contribution and an error no larger than a fixed eventual first term. The shift error is summable for $p\ge0$. Also, for any $\delta\in(0,1)$, the last $(1-\delta)n$ indices have weights $(\log n)^p(1+o(1))$ uniformly, while the first $\delta n$ contribute at most $\delta n(\log n+O(1))^p$. Taking $n\to\infty$ and then $\delta\downarrow0$ gives $W_n\sim n(\log n)^p$. Hence
\[
r_n\sim\frac{\log n}{(p+1)n}.
\]
For the crossing let $N=h^{-1}$ and $g_N=N\log N/(p+1)$. At $n_\pm=\lfloor(1\mp\varepsilon)g_N\rfloor$ and $\lceil(1\pm\varepsilon)g_N\rceil$, respectively, $\log n_\pm/\log N\to1$, so $r_{n_\pm}/h\to(1\mp\varepsilon)^{-1}$. Monotonicity from Appendix~\ref{app:monotonicity} brackets the first crossing between them. Letting $\varepsilon\downarrow0$ gives $n_{{\rm RV},w}\sim g_N$ and, since $n_{\rm pop}\sim N$, the ratio stated in the main text. This includes $p=0$.

\subsection{Log-log weight}\label{app:loglog}
Set $f(x)=\log\log(x+e)>0$. Its derivative is $f'(x)=1/[(x+e)\log(x+e)]$. For large $x$, $f(x)\sim\log\log x$ and $f(x)/x$ is decreasing, so integral comparison applies. The shift from $f(x)$ to $\log\log x$ produces an integrable error after division by $x$: the mean-value theorem bounds it by $C/(x^2\log x)$ for sufficiently large $x$. Substituting $u=\log x$,
\[
A_n=\int^n\frac{\log\log x}{x}\,dx+O(1)
=\log n\log\log n-\log n+O(1)
\sim\log n\log\log n.
\]
The same last-fraction argument as in Appendix~\ref{app:logpower}, using $f(\delta n)/f(n)\to1$ for fixed $\delta>0$, gives $W_n\sim n\log\log n$. Hence $r_n\sim(\log n)/n$. Bracketing at $n_\pm=(1\pm\varepsilon)N\log N$ exactly as in Appendix~\ref{app:logpower} gives $n_{{\rm RV},w}\sim N\log N$ and $n_{{\rm RV},w}/n_{\rm pop}\sim\log N$.

\subsection{Geometric-weight alignment and effective averaging}\label{app:geometric}
For $w_i=q^i$ with $q>1$, write $i=n+1-j$. Division by $W_n$ makes the lag $j$ have probabilities proportional to $q^{-j}$, $j=0,\ldots,n-1$. The denominator $\sum_{j=0}^{n-1}q^{-j}\to q/(q-1)$ and the numerator for $nr_n$ is
\[
\sum_{j=0}^{n-1}q^{-j}\frac{n}{n-j}.
\]
For $j\le n/2$, $n/(n-j)\le2$ and tends to $1$ for each fixed $j$; dominated convergence for the summable geometric series gives the same limit $q/(q-1)$. For $j>n/2$, the sum is at most $n\sum_{j>n/2}q^{-j}\to0$. Thus $nr_n\to1$. Direct finite geometric sums also give
\[
\frac{\sum_{i=2}^{n+1}q^{2i}}{(\sum_{i=2}^{n+1}q^i)^2}
\longrightarrow\frac{q-1}{q+1}.
\]
Theorem~\ref{thm:alignment} supplies the crossing ratio; the displayed square-weight limit measures its bounded effective averaging.

\section{Near-crossing alignment--concentration obstruction}

\subsection{The near-crossing scaling}\label{app:scaling}
Put $h=\theta^2$. Theorem~\ref{thm:population} gives $n_{\rm pop}=h^{-1}+O(1)$. Thus $nh\to c$ places $n$ at asymptotic fraction $c$ of the population crossing. Moreover $n\Delta_R(n)=1-nh\to1-c$, proving the before/after sign interpretation for $c\ne1$.

\subsection{Weight concentration and effective sample size}\label{app:effective}
Let $p_{n,i}=w_i/W_n$, so $\sum_i p_{n,i}=1$. Alignment is
\[
\sum_{i=2}^{n+1}p_{n,i}\left(\frac{n}{i-1}-1\right)\to0.
\]
All summands are nonnegative. For fixed $\varepsilon\in(0,1)$ and $i-1\le(1-\varepsilon)n$, the factor is at least $\varepsilon/(1-\varepsilon)$. Hence the total mass there is at most $(1-\varepsilon)/\varepsilon$ times the displayed sum, and tends to zero. The remaining set has at most $\varepsilon n+1$ indices and mass tending to one. Cauchy--Schwarz on that set gives
\[
\sum_i p_{n,i}^2\ge
\frac{\bigl(\sum_{i-1>(1-\varepsilon)n}p_{n,i}\bigr)^2}
{\varepsilon n+1},
\quad\text{so}\quad
\liminf_n n\sum_i p_{n,i}^2\ge\frac1\varepsilon.
\]
For any $K>0$, choose $\varepsilon<1/K$; the lower limit is then at least $K$. Thus $n\sum_i p_{n,i}^2\to\infty$.

\subsection{Proof of the obstruction theorem}\label{app:obstruction}
Under alignment, $r_n=(1+o(1))/n$. If $nh\to c\ne1$, then by \eqref{eq:weighted-mean}
\[
\frac{\mathbb E\Xi_{n,w}}{W_n}=r_n-h
=\frac{1-c+o(1)}n.
\]
Appendix~\ref{app:full-var} proves that the quadratic variance is nonnegative and the additional linear variance equals $4h\sum_iw_i^2$. Therefore
\[
\frac{\operatorname{Var}(\Xi_{n,w})}{(\mathbb E\Xi_{n,w})^2}
\ge
\frac{4h\sum_i p_{n,i}^2}{(r_n-h)^2}
=
\frac{4(c+o(1))\,n\sum_i p_{n,i}^2}
{(1-c+o(1))^2}\to\infty
\]
by Appendix~\ref{app:effective}. The mean is positive when $c<1$ and negative when $c>1$, for sufficiently large $n$. The proof establishes failure of a sufficient Chebyshev argument, not a positive lower bound on wrong-selection probability. This proves Theorem~\ref{thm:tradeoff}.

\subsection{The exact boundary}\label{app:boundary}
If $nh\to1$ under alignment, $r_n-h=o(1/n)$. When this difference is zero, the variance-to-mean-squared ratio is undefined. When it is nonzero, the preceding lower bound has numerator $4h\sum p_{n,i}^2$ and denominator $o(n^{-2})$; writing $r_n-h=\delta_n/n$ with $\delta_n\to0$ gives a lower bound $4(1+o(1))n\sum p_{n,i}^2/\delta_n^2\to\infty$. No leading-order population sign is assigned at $c=1$.

\section{Exact variance and a fixed-parameter comparison}

\subsection{Loss expansion and centering}\label{app:loss-expansion}
Write $S_m=\sum_{t=1}^m\xi_t$. Since $\bar Y_m=\theta+S_m/m$,
\[
(Y_{m+1}-\bar Y_m)^2-Y_{m+1}^2
=(\xi_{m+1}-S_m/m)^2-(\theta+\xi_{m+1})^2
=\frac{S_m^2}{m^2}-2\xi_{m+1}\frac{S_m}{m}
-2\theta\xi_{m+1}-\theta^2.
\]
The first quadratic term has mean $1/m$ and the other random terms have mean zero. Therefore, with
\[
Q_{n,w}=\sum_{m=1}^nw_{m+1}
\left\{\frac{S_m^2-m}{m^2}-2\xi_{m+1}\frac{S_m}{m}\right\},
\]
we have $\Xi_{n,w}-\mathbb E\Xi_{n,w}=Q_{n,w}-2\theta\sum_{i=2}^{n+1}w_i\xi_i$. This also independently verifies \eqref{eq:weighted-mean}.

\subsection{Gaussian orthogonality and coefficients}\label{app:full-var}
Expand $S_m^2-m=\sum_{t=1}^m(\xi_t^2-1)+2\sum_{1\le t<u\le m}\xi_t\xi_u$. Collecting each centered square gives coefficient
\[
a_t=\sum_{m=t}^n\frac{w_{m+1}}{m^2},\qquad 1\le t\le n.
\]
For each pair $t<l$, the square expansion contributes $2\sum_{m=l}^nw_{m+1}/m^2$; the validation term at $m=l-1$ contributes $-2w_l/(l-1)$. Thus, with empty sums zero,
\[
b_l=\sum_{m=l}^n\frac{w_{m+1}}{m^2}-\frac{w_l}{l-1},
\qquad 2\le l\le n+1,
\]
and
\[
Q_{n,w}=\sum_{t=1}^na_t(\xi_t^2-1)
+2\sum_{l=2}^{n+1}b_l\sum_{t=1}^{l-1}\xi_t\xi_l.
\]
Independent standard Gaussians satisfy $\mathbb E(\xi_t^2-1)^2=2$ and $\mathbb E(\xi_t\xi_l)^2=1$ for $t\ne l$. Distinct centered squares, distinct unordered pair products, and a centered square against a pair product have covariance zero: in every cross-product, some independent centered Gaussian appears to an odd power or as a mean-zero factor. Each quadratic monomial also has zero covariance with every $\xi_i$, because every centered third Gaussian moment is zero. Hence the quadratic and linear parts are orthogonal, and
\[
\operatorname{Var}(Q_{n,w})=2\sum_{t=1}^na_t^2+
4\sum_{l=2}^{n+1}(l-1)b_l^2,\qquad
\operatorname{Var}\!\left(-2\theta\sum_{i=2}^{n+1}w_i\xi_i\right)
=4\theta^2\sum_{i=2}^{n+1}w_i^2.
\]
Adding proves the exact identity \eqref{eq:variance}.

\subsection{Slow-weight variance orders}\label{app:slow-var}
Here $h=\theta^2$ and $nh\to c\in(0,\infty)$; consider $w_{m+1}=f(m)$, where $f(m)=1$, $[\log(m+2)]^p$ for fixed $p>0$, or $\log\log(m+1+e)$. In each case $f(n)$ is slowly varying, and splitting sums at $\delta n$ as in Appendix B shows
\[
\sum_{m=1}^n f(m)^2\sim nf(n)^2.
\]
We verify the quadratic order without borrowing it from 08. The coefficients $a_t$ are bounded by $\sum_{m=t}^\infty f(m)/m^2$. For the listed $f$, integral comparison and eventual monotonicity of $f(x)/x^{3/4}$ give $\sum_{m=t}^\infty f(m)/m^2\le C f(t)/t$ for large $t$. Consequently $\sum_t a_t^2\le C+C\sum_t f(t)^2/t^2=O(1)$.

Set $T_n=\sum_{m=n+1}^\infty f(m)/m^2$ and $B_l=\sum_{m=l}^\infty f(m)/m^2-f(l-1)/(l-1)$, so $b_l=B_l-T_n$ for $2\le l\le n$ and $b_{n+1}=-f(n)/n$ (the harmless one-step indexing changes the following estimates only by lower-order terms). Integral comparison, summation by parts, and the mean-value theorem give
\[
T_n\sim f(n)/n,\qquad
|B_l|\le C\left(\frac{\sup_{x\in[l-1,l+1]}|f'(x)|}{1}
+\frac{f(l)}{l^2}\right)
\]
for smooth extensions of the listed $f$. To see the derivative bound, subtract the constant-$f(l)$ sum, whose discrepancy from $f(l)/(l-1)$ is $O(f(l)/l^2)$; summation by parts bounds the remaining variation by $C\int_l^\infty |f'(x)|x^{-1}dx$. This is $O((\log l)^{p-1}/l)$ for log powers and $O(1/(l\log l))$ for log-log; the constant case has zero derivative. Therefore $D_n=\sum_{l=2}^n(l-1)B_l^2$ is $O(1)$ for $f=1$ and log-log, while for log powers it is bounded by
\[
C+C\sum_{l=3}^n\frac{(\log l)^{2p-2}}l
=\begin{cases}
O(1),&0<p<1/2,\\
O(\log\log n),&p=1/2,\\
O((\log n)^{2p-1}),&p>1/2.
\end{cases}
\]
The three cases follow by $u=\log x$ in the comparison integral $\int^n(\log x)^{2p-2}dx/x$. In every case $D_n=o(f(n)^2)$ when $f(n)\to\infty$; for $f=1$, $D_n=O(1)$ and the exact $B_l=O(l^{-2})$ is handled below. The terminal contribution is
\[
4T_n^2\sum_{l=2}^n(l-1)\sim
4\frac{f(n)^2}{n^2}\frac{n^2}{2}=2f(n)^2.
\]
Expanding $(B_l-T_n)^2$, Cauchy--Schwarz bounds the magnitude of its cross term by
$2\{D_n T_n^2\sum_{l=2}^n(l-1)\}^{1/2}=o(f(n)^2)$ for the growing weights. The end coefficient contributes $4n(f(n)/n)^2=o(f(n)^2)$. Thus $\operatorname{Var}(Q_{n,w})\sim2f(n)^2$ for the growing log and log-log weights. For $f=1$, direct integral comparison yields $B_l=O(l^{-2})$. Thus $D_n=O(1)$ and $\sum_{l=2}^n(l-1)|B_l|=O(\log n)$. Since $T_n=O(1/n)$, the cross term is $O(\log n/n)=o(1)$; the end coefficient is $O(1/n)$. Therefore $\operatorname{Var}(Q_{n,w})\sim2$ also in the unweighted case.

The linear variance is $4h\sum_iw_i^2\sim4c f(n)^2$. Hence the full variance is of order $f(n)^2$ (indeed asymptotic to $(2+4c)f(n)^2$ for the growing weights). Meanwhile Appendix B gives $\mathbb E\Xi_{n,w}\sim(\log n)f(n)/(p+1)$ for log powers and $\sim(\log n)f(n)$ for log-log; the $-hW_n$ term is only $O(f(n))$. The ordinary-weight mean is $\sim\log n$ and its full variance tends to $2+4c$. Therefore all four Chebyshev ratios have order $(\log n)^{-2}$. Their means are positive for fixed $c$, establishing the sign statement in the main text.

\subsection{Stretched-exponential weights at fixed $\theta$}\label{app:stretched}
Let $w_i=\exp(i^\beta)$ with $0<\beta<1$. Put $F(x)=e^{x^\beta}$. Substitution $u=x^\beta$ and integration by parts yield, for $a=1,2$,
\[
\int^n F(x)^a\,dx
=\frac1\beta\int^{n^\beta} e^{au}u^{1/\beta-1}\,du
\sim\frac{n^{1-\beta}e^{an^\beta}}{a\beta}.
\]
Because $F$ is increasing, bounding sums between adjacent integrals changes them by at most $O(F(n)^a)$, negligible relative to $n^{1-\beta}F(n)^a$. Endpoint shifts by one are negligible since $(n+1)^\beta-n^\beta\to0$. Therefore
\[
W_n\sim\frac{n^{1-\beta}e^{n^\beta}}{\beta},\quad
\sum_iw_i^2\sim\frac{n^{1-\beta}e^{2n^\beta}}{2\beta},\quad
\frac{\sum_iw_i^2}{W_n^2}\sim\frac{\beta}{2n^{1-\beta}}\to0.
\]
These sums show an effective terminal window of order $n^{1-\beta}=o(n)$. For any fixed $\delta>0$, the total weight at $i-1\le(1-\delta)n$ is at most $n e^{((1-\delta)n+O(1))^\beta}=o(W_n)$, because the exponential gap is of order $n^\beta$. On the remaining indices, $n/(i-1)\le1/(1-\delta)$; for the earlier indices their weighted contribution to $nr_n$ is negligible by the same exponential gap even after the factor $n/(i-1)\le n$. Hence $1\le\liminf nr_n\le\limsup nr_n\le1/(1-\delta)$, and $\delta\downarrow0$ proves alignment.

For fixed $\theta\ne0$, $\mathbb E\Xi_{n,w}/W_n=r_n-\theta^2\to-\theta^2$. To bound the quadratic variance, define the centered single-time quadratic increment $U_m=(S_m^2-m)/m^2-2\xi_{m+1}S_m/m$. Gaussian moments and independence give
$\|U_m\|_2\le\sqrt2/m+2/\sqrt m\le C/\sqrt m$.
Minkowski's inequality and then Cauchy--Schwarz yield
\[
\left\|\frac{Q_{n,w}}{W_n}\right\|_2
\le C\sum_{m=1}^n\frac{w_{m+1}}{W_n\sqrt m}
\le C\left(\sum_{m=1}^n\frac{w_{m+1}}{W_nm}\right)^{1/2}
=C\sqrt{r_n}\to0.
\]
The normalized linear variance is $4\theta^2\sum_iw_i^2/W_n^2\to0$. Orthogonality from Appendix~\ref{app:full-var} then gives $\operatorname{Var}(\Xi_{n,w})/W_n^2\to0$, and division by the mean-squared limit $\theta^4$ proves fixed-$\theta$ concentration. This does not apply to $\theta\to0$ with $n\theta^2\to c$.

\end{document}